\chapter{Fundamente teoretice si tehnologice}

\section{Paradigme si principii de proiectare}

Trainee urmeaza un model client-server clasic, dar bine delimitat: interfata mobila este separata de logica de business si de persistenta datelor. Aceasta separare simplifica testarea, mentenanta si, foarte important, controlul asupra punctelor sensibile de securitate.

In proiectare au fost urmate principii standard de inginerie software:
\begin{itemize}[leftmargin=1.2cm]
    \item \textbf{Single Responsibility}: fiecare componenta trateaza un set restrans de responsabilitati;
    \item \textbf{Separation of Concerns}: validarea, autentificarea, autorizarea si limitarea traficului sunt tratate in middleware-uri distincte;
    \item \textbf{Fail Fast}: configuratiile critice sunt validate inainte ca aplicatia sa serveasca trafic;
    \item \textbf{Defence in Depth}: fluxurile sensibile sunt protejate pe mai multe niveluri, nu printr-un singur control.
\end{itemize}

Aceste principii raman coerente cu recomandarile generale din literatura de specialitate \parencite{pressman2014software}.

\section{Tehnologii utilizate}

\subsection{Frontend mobil}

Partea client este dezvoltata cu Expo si React Native, iar navigatia este gestionata cu Expo Router. Pentru stare si comunicare API este folosit Redux Toolkit impreuna cu RTK Query \parencite{reduxdocs}. Combinatia a fost aleasa pentru ca ofera un echilibru bun intre viteza de dezvoltare si predictibilitatea comportamentului.

Avantaje practice observate:
\begin{itemize}[leftmargin=1.2cm]
    \item iteratie rapida pe dispozitive mobile;
    \item integrare naturala cu ecosistemul React;
    \item gestionare unitara pentru cache, refetch si actualizari de stare;
    \item reducerea codului repetitiv in ecranele care consuma endpoint-uri similare.
\end{itemize}

\subsection{Backend API}

Backend-ul este construit cu Node.js, Express si TypeScript. Express a fost preferat datorita flexibilitatii si usurintei de integrare cu middleware-uri custom \parencite{expressdocs}. TypeScript adauga un strat important de siguranta prin tipuri explicite pe payload-uri, modele si raspunsuri.

\subsection{Persistenta si cautare geospatiala}

Persistenta foloseste PostgreSQL prin Sequelize, cu extensii PostGIS pentru proximitate si pg\_trgm pentru cautare text optimizata \parencite{postgresdocs,postgisdocs}. In termeni functionali, asta inseamna rezultate de cautare mai relevante in functie de locatie si performanta stabila la cresterea datelor.

\subsection{Servicii externe}

Platforma se integreaza cu servicii externe pentru functionalitati specializate:
\begin{itemize}[leftmargin=1.2cm]
    \item Stripe si RevenueCat pentru billing si entitlement;
    \item SMTP/Nodemailer pentru emailuri de sistem;
    \item AWS S3 pentru stocare media;
    \item surse externe pentru popularea datelor initiale de sali.
\end{itemize}

\section{Securitate aplicata}

Abordarea de securitate este inspirata din recomandarile OWASP Top 10 \parencite{owasp2021top10}, dar aplicata pragmatic in cod:
\begin{itemize}[leftmargin=1.2cm]
    \item autentificare JWT pentru endpoint-uri protejate;
    \item validare stricta pe \texttt{body}, \texttt{query} si \texttt{params};
    \item profile de rate limiting in functie de tipul endpoint-ului;
    \item autorizare pe rol pentru operatiuni administrative;
    \item gestiune explicita a secretelor prin variabile de mediu.
\end{itemize}

Pe langa aceste mecanisme, proiectul include reguli anti-abuz pentru vizualizarile de profil public, astfel incat indicatorii de analytics sa nu poata fi manipulati usor.

\section{Alternative analizate}

Inainte de alegerea stack-ului final au fost evaluate mai multe optiuni.

\begin{table}[H]
\centering
\caption{Alternative tehnologice evaluate}
\begin{tabular}{p{0.23\textwidth} p{0.27\textwidth} p{0.40\textwidth}}
\toprule
\textbf{Domeniu} & \textbf{Optiuni evaluate} & \textbf{Argument pentru alegerea finala} \\
\midrule
Frontend mobil & React Native / Flutter & React Native + Expo a oferit timp de livrare mai bun si integrare mai directa cu expertiza existenta. \\
Backend API & Express / NestJS / Fastify & Express a oferit libertate maxima pentru un design incremental al middleware-urilor. \\
ORM & Sequelize / Prisma / TypeORM & Sequelize era deja compatibil cu structura proiectului si modelele existente. \\
Persistenta & PostgreSQL / MongoDB & Cerintele relationale si suportul geospatial au favorizat PostgreSQL. \\
Billing & Stripe-only / RevenueCat-only / hibrid & Modelul hibrid permite tranzitie operationala fara rescrieri majore de flux. \\
\bottomrule
\end{tabular}
\end{table}

\section{Cadru academic si conformitate}

Redactarea lucrarii urmeaza reperele publicate de FMI pentru finalizarea studiilor \parencite{fmi2026finalizare}. Folosirea LaTeX/Overleaf ofera control bun asupra structurii, numerotarii, citarilor si anexelor tehnice.

In acest cadru, capitolele urmatoare detaliaza modul in care deciziile tehnologice sunt transformate in implementare concreta.
